% GENERATED FILE. Edit canonical lessons, then run npm run generate:books.
% canonical-source-hash: fnv1a64:65fb6f59550454fb
% canonical-lessons: BN-C118-rosogolla, BN-C118-payes, BN-C118-luchi, BN-C118-muri, BN-C118-kophi

\chapter{Rasgulla, Rice pudding, Luchi, Puffed rice, Coffee}
\label{ch:bn-rasgulla-rice-pudding-luchi-puffed-rice-coffee}
\hlchaptermodality{118}

\begin{chapteropening}
\textbf{By the end of this chapter:} \emph{I can say rasgulla, rice pudding, a luchi, puffed rice and coffee.}

Complete the last lesson of chapter 118: I can say rasgulla, rice pudding, a luchi, puffed rice and coffee.
\end{chapteropening}

\section[\texorpdfstring{\bn{রসগোল্লা} (rôsôgollā)}{রসগোল্লা (rôsôgollā)}]{\bn{রসগোল্লা} (rôsôgollā) — rasgulla}

\label{lesson:BN-C118-rosogolla}



\begin{quote}
Before the new one: say the Bengali for a radish, then the Bengali for uncooked rice.
\end{quote}

\subsection*{You'll want to know: \bn{রসগোল্লা}}
\textbf{\bn{রসগোল্লা}} — \emph{rôsôgollā} — \textquotedblleft{}rasgulla\textquotedblright{}.

\begin{grammarlens}[title={What you've built}]
One more word for this chapter.
\end{grammarlens}

\subsection*{Your turn}
\begin{itemize}
\raggedright
  \item \emph{Say it:} \emph{rôsôgollā}
  \item \emph{Say it:} \emph{rôsôgollā}, once more
  \item \emph{Say it:} say \emph{rôsôgollā}
  \item \emph{Recall:} say the Bengali for a radish, then the Bengali for uncooked rice, then say \emph{rôsôgollā} again
\end{itemize}

\begin{tcolorbox}[breakable,colback=teal!4,colframe=teal!35!black,title={Before you move on}]
What does \bn{রসগোল্লা} mean? (\textquotedblleft{}Rasgulla\textquotedblright{}.) Say it once more.
\end{tcolorbox}
\section[\texorpdfstring{\bn{পায়েস} (pāyes)}{পায়েস (pāyes)}]{\bn{পায়েস} (pāyes) — rice pudding}

\label{lesson:BN-C118-payes}



\begin{quote}
Before the new one: say the Bengali for uncooked rice, then the Bengali for rasgulla.
\end{quote}

\subsection*{You'll want to know: \bn{পায়েস}}
\textbf{\bn{পায়েস}} — \emph{pāyes} — \textquotedblleft{}rice pudding\textquotedblright{}.

\begin{grammarlens}[title={What you've built}]
One more word for this chapter.
\end{grammarlens}

\subsection*{Your turn}
\begin{itemize}
\raggedright
  \item \emph{Say it:} \emph{pāyes}
  \item \emph{Say it:} \emph{pāyes}, once more
  \item \emph{Say it:} say \emph{pāyes}
  \item \emph{Recall:} say the Bengali for uncooked rice, then the Bengali for rasgulla, then say \emph{pāyes} again
\end{itemize}

\begin{tcolorbox}[breakable,colback=teal!4,colframe=teal!35!black,title={Before you move on}]
What does \bn{পায়েস} mean? (\textquotedblleft{}Rice pudding\textquotedblright{}.) Say it once more.
\end{tcolorbox}
\section[\texorpdfstring{\bn{লুচি} (luchi)}{লুচি (luchi)}]{\bn{লুচি} (luchi) — a luchi}

\label{lesson:BN-C118-luchi}



\begin{quote}
Before the new one: say the Bengali for rasgulla, then the Bengali for rice pudding.
\end{quote}

\subsection*{You'll want to know: \bn{লুচি}}
\textbf{\bn{লুচি}} — \emph{luchi} — \textquotedblleft{}a luchi\textquotedblright{}.

\begin{grammarlens}[title={What you've built}]
One more word for this chapter.
\end{grammarlens}

\subsection*{Your turn}
\begin{itemize}
\raggedright
  \item \emph{Say it:} \emph{luchi}
  \item \emph{Say it:} \emph{luchi}, once more
  \item \emph{Say it:} say \emph{luchi}
  \item \emph{Recall:} say the Bengali for rasgulla, then the Bengali for rice pudding, then say \emph{luchi} again
\end{itemize}

\begin{tcolorbox}[breakable,colback=teal!4,colframe=teal!35!black,title={Before you move on}]
What does \bn{লুচি} mean? (\textquotedblleft{}A luchi\textquotedblright{}.) Say it once more.
\end{tcolorbox}
\section[\texorpdfstring{\bn{মুড়ি} (muṛi)}{মুড়ি (muṛi)}]{\bn{মুড়ি} (muṛi) — puffed rice}

\label{lesson:BN-C118-muri}



\begin{quote}
Before the new one: say the Bengali for rice pudding, then the Bengali for a luchi.
\end{quote}

\subsection*{You'll want to know: \bn{মুড়ি}}
\textbf{\bn{মুড়ি}} — \emph{muṛi} — \textquotedblleft{}puffed rice\textquotedblright{}.

\begin{grammarlens}[title={What you've built}]
One more word for this chapter.
\end{grammarlens}

\subsection*{Your turn}
\begin{itemize}
\raggedright
  \item \emph{Say it:} \emph{muṛi}
  \item \emph{Say it:} \emph{muṛi}, once more
  \item \emph{Say it:} say \emph{muṛi}
  \item \emph{Recall:} say the Bengali for rice pudding, then the Bengali for a luchi, then say \emph{muṛi} again
\end{itemize}

\begin{tcolorbox}[breakable,colback=teal!4,colframe=teal!35!black,title={Before you move on}]
What does \bn{মুড়ি} mean? (\textquotedblleft{}Puffed rice\textquotedblright{}.) Say it once more.
\end{tcolorbox}
\section[\texorpdfstring{\bn{কফি} (kôphi)}{কফি (kôphi)}]{\bn{কফি} (kôphi) — coffee}

\label{lesson:BN-C118-kophi}



\begin{quote}
Before the new one: say the Bengali for a luchi, then the Bengali for puffed rice.
\end{quote}

\subsection*{You'll want to know: \bn{কফি}}
\textbf{\bn{কফি}} — \emph{kôphi} — \textquotedblleft{}coffee\textquotedblright{}.

\begin{grammarlens}[title={What you've built}]
One more word for this chapter.
\end{grammarlens}

\subsection*{Your turn}
\begin{itemize}
\raggedright
  \item \emph{Say it:} \emph{kôphi}
  \item \emph{Say it:} \emph{kôphi}, once more
  \item \emph{Say it:} say \emph{kôphi}
  \item \emph{Recall:} say the Bengali for a luchi, then the Bengali for puffed rice, then say \emph{kôphi} again
\end{itemize}

\begin{tcolorbox}[breakable,colback=teal!4,colframe=teal!35!black,title={Before you move on}]
What does \bn{কফি} mean? (\textquotedblleft{}Coffee\textquotedblright{}.) Say it once more.
\end{tcolorbox}
